\documentclass[runningheads]{llncs}

% ============================================================
% Packages
% ============================================================
\usepackage{cite}
\usepackage{amsmath,amssymb,amsfonts}
\usepackage{algorithmic}
\usepackage{algorithm}
\usepackage{graphicx}
\usepackage{textcomp}
\usepackage{xcolor}
\usepackage{booktabs}
\usepackage{multirow}
\usepackage{tabularx}
\usepackage{array}
% 'float' package helps fix timeout issues by forcing placement
\usepackage{float} 

% ============================================================
% Space Saving & Timeout Prevention Settings
% ============================================================
% 1. Reduce space around tables and figures
\setlength{\textfloatsep}{5pt plus 1.0pt minus 2.0pt}
\setlength{\floatsep}{5pt plus 1.0pt minus 2.0pt}
\setlength{\intextsep}{5pt plus 1.0pt minus 2.0pt}

% 2. Reduce space around equations
\AtBeginDocument{
  \setlength{\abovedisplayskip}{4pt}
  \setlength{\belowdisplayskip}{4pt}
  \setlength{\abovedisplayshortskip}{0pt}
  \setlength{\belowdisplayshortskip}{0pt}
}

% 3. Tighten bibliography spacing (Define if not present in llncs, though normally strict)
% \def\IEEEbibitemsep{0pt plus .5pt} % Removed as it is IEEE specific.

% ============================================================
\begin{document}

\title{DetSEG: A Two-Stage Coarse-to-Fine Framework for Automated Polyp Detection and Segmentation in Colonoscopy Images}

\author{Author Name\inst{1}}
\institute{Department of Computer Science, Institution Name, City, Country\\
\email{email@institution.edu}}

\maketitle

% ============================================================
% ABSTRACT
% ============================================================
\begin{abstract}
Colorectal cancer (CRC) remains one of the leading causes of cancer-related mortality worldwide. Early detection and removal of colonic polyps during colonoscopy is the most effective means of preventing CRC. However, polyp detection during routine colonoscopy is highly dependent on operator skill, with miss rates ranging from 6\% to 27\%. This paper presents DetSEG, a novel two-stage coarse-to-fine deep learning framework for automated polyp detection and segmentation in colonoscopy images. The proposed architecture employs YOLOv8 for real-time polyp localization and a U-Net with ResNet-34 backbone for pixel-wise segmentation of detected regions. Our ``detect-then-segment'' approach addresses the challenge of polyp scale variation and background complexity by decoupling the localization and segmentation tasks, ensuring focused feature extraction. Experimental results demonstrate that DetSEG achieves superior segmentation accuracy compared to end-to-end approaches while maintaining computational efficiency suitable for clinical deployment. The framework achieves a mean Dice coefficient of 0.89 and Intersection over Union (IoU) of 0.84 on benchmark datasets, with inference times of under 100ms per image.

\keywords{Polyp Detection \and Semantic Segmentation \and Colonoscopy \and Deep Learning \and U-Net \and YOLO \and Computer-Aided Diagnosis}
\end{abstract}

% ============================================================
% 1. INTRODUCTION
% ============================================================
\section{Introduction}

\subsection{Background and Motivation}

Colorectal cancer (CRC) is the third most commonly diagnosed malignancy and the second leading cause of cancer-related deaths globally, accounting for approximately 1.9 million new cases and 935,000 deaths annually. The adenoma-carcinoma sequence, wherein benign adenomatous polyps progressively transform into malignant carcinomas over a period of 10--15 years, provides a critical window for early intervention through colonoscopic polypectomy. Colonoscopy serves as the gold standard for colorectal cancer screening, enabling both detection and removal of precancerous polyps. However, the adenoma detection rate (ADR)---defined as the proportion of screening colonoscopies during which at least one adenomatous polyp is detected---varies significantly among endoscopists. Studies indicate polyp miss rates of 6--27\% during routine colonoscopy, influenced by factors including:

\noindent Operator fatigue : Prolonged examination times lead to decreased vigilance

\noindent Polyp morphology: Flat and sessile polyps are inherently more difficult to detect

\noindent Bowel preparation quality: Suboptimal bowel cleansing obscures mucosal visualization

\noindent Optical limitations: Blind spots in colonic folds and angulated segments

The critical need to reduce ADR variability and minimize polyp miss rates has catalyzed research into computer-aided detection (CADe) and computer-aided diagnosis (CADx) systems for colonoscopy. Deep learning approaches have shown remarkable promise in this domain, offering the potential for real-time, consistent, and objective polyp detection that can complement endoscopist expertise.

Despite significant advancements, automated polyp analysis remains challenging due to several inherent complexities:

\noindent Scale Variation: Polyps vary dramatically in size, from diminutive (5mm) to large ($>$20mm), requiring multi-scale feature extraction

\noindent Morphological Diversity: The Paris classification defines multiple polyp morphologies (pedunculated, sessile, flat, depressed) with distinct visual characteristics

\noindent Color and Texture Similarity: Polyps often exhibit color and texture similar to surrounding normal mucosa

\noindent Illumination Variability: Non-uniform lighting during colonoscopy creates specular reflections and shadows

\noindent Motion Artifacts: Patient and scope movement introduces blur and distortion

\noindent Background Complexity: Colonic folds, vessels, and residual matter create cluttered backgrounds

These challenges motivate the development of robust frameworks that can effectively handle the variability and complexity inherent in colonoscopy images.

\subsection{Contributions}

This paper presents DetSEG, a two-stage deep learning framework for polyp detection and segmentation. The main contributions of this work are:

\noindent Novel Architecture Design: We propose a coarse-to-fine framework that decouples object detection from semantic segmentation, enabling each module to specialize in its respective task

\noindent Optimized Detection Pipeline: We adapt YOLOv8 for polyp localization, leveraging its anchor-free detection paradigm for improved generalization on irregular polyp shapes

\noindent Transfer Learning Strategy: We utilize ImageNet-pretrained ResNet-34 as the U-Net encoder, enabling effective feature extraction even with limited medical imaging data

\noindent Clinical Applicability: The system achieves sub-100ms inference times suitable for clinical deployment

\noindent Comprehensive Evaluation: We provide extensive ablation studies and comparison with state-of-the-art methods on benchmark datasets

% ============================================================
% 2. RELATED WORK
% ============================================================
\section{Related Work}

\subsection{Traditional Computer-Aided Detection Methods}

Early CADe systems for polyp detection relied on hand-crafted features and classical machine learning algorithms. Notable approaches include:

\noindent Texture-based methods: Utilizing local binary patterns (LBP), Gabor filters, and co-occurrence matrices to characterize polyp surface texture

\noindent Shape-based methods: Employing curvature analysis, ellipse fitting, and protrusion detection to identify polypoid morphology

\noindent Color-based methods: Exploiting color histogram analysis in various color spaces (RGB, HSV, Lab) to distinguish polyps from normal mucosa

\subsection{Deep Learning for Polyp Detection}

The advent of convolutional neural networks (CNNs) revolutionized medical image analysis. Object detection architectures that have been applied to polyp detection include:

\subsubsection{Two-Stage Detectors}
R-CNN Family: Girshick et al. introduced R-CNN, later refined through Fast R-CNN and Faster R-CNN. These methods generate region proposals followed by classification and bounding box regression.

\subsubsection{Single-Stage Detectors}
YOLO (You Only Look Once): Redmon et al. pioneered real-time object detection with YOLO, treating detection as a regression problem. YOLOv8 represents the latest evolution with improved accuracy and speed through:

\noindent C2f (Cross-Stage Partial bottleneck with 2 convolutions and flow) modules

\noindent Anchor-free split Ultralytics head

\noindent Distribution Focal Loss for improved localization

SSD (Single Shot Detector): Multi-scale feature maps for detecting objects at various sizes.

RetinaNet: Focal loss to address class imbalance in object detection.

\subsection{Deep Learning for Polyp Segmentation}

Semantic segmentation networks provide pixel-wise labeling, crucial for accurate polyp boundary delineation.

\subsubsection{U-Net and Variants}
U-Net, introduced by Ronneberger et al. for biomedical image segmentation, features:

\noindent Encoder-decoder architecture with skip connections

\noindent Precise localization through concatenation of high-resolution encoder features

\noindent Strong performance with limited training data through data augmentation

Variants include U-Net++, Attention U-Net, and ResUNet, which incorporate nested skip pathways, attention mechanisms, and residual connections, respectively.

\subsubsection{Polyp-Specific Architectures}
\noindent PraNet: Parallel Reverse Attention Network with area attention modules for polyp segmentation

\noindent SANet: Self-Attention Network with multi-scale context aggregation

\noindent HarDNet-MSEG: High-resolution deep network for medical image segmentation

\noindent Polyp-PVT: Pyramid Vision Transformer for polyp segmentation

\subsection{Two-Stage Detection-Segmentation Frameworks}

Several works have explored cascaded detection-segmentation pipelines:

\noindent Mask R-CNN: Extends Faster R-CNN with a branch for predicting segmentation masks within each Region of Interest (RoI)

\noindent YOLACT/YOLACT++: Real-time instance segmentation through prototype-based mask generation

% ============================================================
% 3. PROPOSED METHODOLOGY
% ============================================================
\section{Proposed Methodology}

\subsection{DetSEG Framework Overview}

The DetSEG framework employs a coarse-to-fine strategy comprising two sequential modules:

\noindent 1. Polyp Localization Module (PLM): A YOLOv8-based detector that identifies candidate polyp regions

\noindent 2. Fine-Grained Segmentation Module (FSM): A U-Net with ResNet-34 encoder that generates pixel-wise masks for detected regions

\begin{figure}[htbp]
\centering
%\includegraphics[width=\textwidth]{model_architecture.png}
\caption{The proposed DetSEG framework architecture. The first stage (Polyp Localization Module) uses YOLOv8 to detect polyp bounding boxes. The second stage (Fine-Grained Segmentation Module) uses a U-Net with ResNet-34 encoder to segment the polyp within the cropped bounding box.}
\label{fig:architecture}
\end{figure}

This architecture offers several advantages:

\noindent Reduced Search Space: The segmentation network focuses exclusively on detected regions, minimizing false positives from background structures

\noindent Scale Normalization: Resizing detected regions to fixed dimensions addresses polyp scale variation

\noindent Modular Design: Each component can be independently trained and optimized

\noindent Computational Efficiency: Only regions of interest undergo the more expensive segmentation process

% 

\subsection{Dataset Preparation and Preprocessing}

\subsubsection{Dataset Description}
The model development utilizes publicly available colonoscopy datasets containing RGB images with corresponding binary segmentation masks:

\noindent Kvasir-SEG: 1,000 polyp images with ground truth masks

\noindent CVC-ClinicDB: 612 frames from 31 colonoscopy sequences

\subsubsection{Annotation Transformation}
To support the two-stage architecture, we transformed segmentation annotations into detection annotations. Bounding Box Generation: Minimum bounding rectangles were computed from binary ground-truth masks. For a mask M with polyp pixels $P=\{(x_{i},y_{i}) : M(x_{i},y_{i})=1\}$, the bounding box coordinates are:

\begin{equation}
x_{min}=min_{(x_{i},y_{i})\in P}x_{i}, \quad y_{min}=min_{(x_{i},y_{i})\in P}y_{i}
\end{equation}
\begin{equation}
x_{max}=max_{(x_{i},y_{i})\in P}x_{i}, \quad y_{max}=max_{(x_{i},y_{i})\in P}y_{i}
\end{equation}

The bounding box is then represented in YOLO format as normalized center coordinates and dimensions:

\begin{equation}
x_{c}=\frac{x_{min}+x_{max}}{2W}, \quad y_{c}=\frac{y_{min}+y_{max}}{2H}
\end{equation}
\begin{equation}
w_{norm}=\frac{x_{max}-x_{min}}{W}, \quad h_{norm}=\frac{y_{max}-y_{min}}{H}
\end{equation}

where W and H are the image width and height, respectively.

\subsubsection{Patch Generation for Segmentation Training}
For the FSM, we generated a dataset of polyp patches:

\noindent 1. Crop Extraction: Using ground-truth bounding boxes, polyp regions were extracted with contextual padding ($p=20$ pixels) to provide spatial context

\noindent 2. Padding Calculation:
\begin{equation}
x_{1}^{\prime}=max(0,x_{min}-p), \quad y_{1}^{\prime}=max(0,y_{min}-p)
\end{equation}
\begin{equation}
x_{2}^{\prime}=min(W,x_{max}+p), \quad y_{2}^{\prime}=min(H,y_{max}+p)
\end{equation}

\noindent 3. Resizing: Patches were resized to $256\times256$ pixels using bilinear interpolation for images and nearest-neighbor interpolation for masks to preserve binary values

\subsubsection{Data Split}
The dataset was partitioned with stratified random sampling (seed=42).

\begin{table}[H]
\caption{Data Split}
\label{tab:data_split}
\centering
\begin{tabular}{lcl}
\toprule
\textbf{Split} & \textbf{Proportion} & \textbf{Purpose} \\
\midrule
Train & 80\% & Model training \\
Validation & 10\% & Hyperparameter tuning \\
Test & 10\% & Final evaluation \\
\bottomrule
\end{tabular}
\end{table}

\subsection{Stage 1: Polyp Localization Module (PLM)}

\subsubsection{YOLOv8 Architecture}
We employed YOLOv8s (small variant), selected for its optimal balance between detection accuracy and inference speed. The architecture comprises:

\noindent Backbone (CSPDarknet53): Cross-Stage Partial Darknet with 53 layers. Efficiently processes input through hierarchical feature extraction. Progressive spatial reduction with feature dimension increase.

\noindent Neck (FPN + PANet): Feature Pyramid Network (FPN) for top-down feature propagation. Path Aggregation Network (PANet) for bottom-up pathway. Multi-scale feature fusion for detecting polyps of varying sizes.

\noindent Head (Anchor-Free): Decoupled head for classification and regression. Anchor-free design eliminates manual anchor configuration. Distribution Focal Loss for precise bounding box prediction.

The architecture can be formalized as:
\begin{equation}
F_{backbone}=CSPDarknet(I_{input})
\end{equation}
\begin{equation}
F_{neck}=PANet(FPN(F_{backbone}))
\end{equation}
\begin{equation}
\{B,C,P\}=Head(F_{neck})
\end{equation}
where B represents bounding boxes, C represents class predictions, and P represents confidence scores.

\subsubsection{YOLOv8 Loss Function}
The total loss combines three components:

\noindent Box Regression Loss (CIoU): The Complete IoU loss addresses limitations of standard IoU by considering overlap, center distance, and aspect ratio:

\begin{equation}
\mathcal{L}_{CIoU}=1-IoU+\frac{\rho^{2}(b,b^{gt})}{c^{2}}+\alpha v
\end{equation}

where: $IoU=\frac{|B \cap B^{gt}|}{|B \cup B^{gt}|}$ is the standard Intersection over Union; $\rho(b,b^{gt})$ is the Euclidean distance between predicted and ground truth box centers; c is the diagonal length of the smallest enclosing box; and v encode aspect ratio consistency:

\begin{equation}
v=\frac{4}{\pi^{2}}(arctan\frac{w^{gt}}{h^{gt}}-arctan\frac{w}{h})^{2}
\end{equation}
\begin{equation}
\alpha=\frac{v}{(1-IoU)+v}
\end{equation}

\noindent Classification Loss (Binary Cross-Entropy):

\begin{equation}
\mathcal{L}_{cls}=-\sum_{i=1}^{N}[y_{i}log(\tilde{y}_{i})+(1-y_{i})log(1-\tilde{y}_{i})]
\end{equation}

\noindent Distribution Focal Loss (DFL): For precise localization, DFL models bounding box boundaries as probability distributions:

\begin{equation}
\mathcal{L}_{DFL}=-((y_{i+1}-y)log(S_{i})+(y-y_{i})log(S_{i+1}))
\end{equation}

where y lies between $y_{i}$ and $y_{i+1}$ and S represents the softmax output.

\noindent Total Detection Loss:

\begin{equation}
\mathcal{L}_{det}=\lambda_{1}\mathcal{L}_{CIoU}+\lambda_{2}\mathcal{L}_{cls}+\lambda_{3}\mathcal{L}_{DFL}
\end{equation}

\subsubsection{PLM Training Configuration}
\begin{table}[H]
\caption{PLM Training Configuration}
\label{tab:plm_config}
\centering
\begin{tabular}{ll}
\toprule
\textbf{Hyperparameter} & \textbf{Value} \\
\midrule
Input Resolution & $640\times640$ \\
Batch Size & 16 \\
Epochs & 100 \\
Optimizer & SGD with momentum (0.937) \\
Initial Learning Rate & 0.01 \\
LR Schedule & Cosine annealing \\
Weight Decay & 0.0005 \\
Pretrained Weights & COCO \\
Data Augmentation & Mosaic, HSV shift, Flip \\
\bottomrule
\end{tabular}
\end{table}

\subsubsection{Non-Maximum Suppression (NMS)}
During inference, overlapping detections are refined using NMS.

\begin{algorithm}[H]
\caption{Non-Maximum Suppression (NMS)}
\label{alg:nms}
\begin{algorithmic}[1]
\STATE Sort detections by confidence score in descending order
\STATE Select the detection with highest confidence
\STATE Compute IoU with all remaining detections
\STATE Suppress detections with IoU $>$ threshold (0.45)
\STATE Repeat until no detections remain
\end{algorithmic}
\end{algorithm}

\subsection{Stage 2: Fine-Grained Segmentation Module (FSM)}

\subsubsection{U-Net Architecture with ResNet-34 Encoder}
The FSM employs a U-Net architecture enhanced with a pretrained ResNet-34 encoder.


\noindent Encoder (Contracting Path): The ResNet-34 encoder replaces the traditional U-Net contracting path:

\noindent Layer 0 (Initial): $7\times7$ convolution, stride 2, 64 filters, followed by batch normalization, ReLU, and $3\times3$ max pooling with stride 2

\noindent Layers 1-4: Residual blocks with increasing channel dimensions [64, 128, 256, 512]

The residual connection is defined as:
\begin{equation}
y=\mathcal{F}(x,\{W_{i}\})+x
\end{equation}
where $\mathcal{F}$ represents the learned residual function.

\noindent Decoder (Expanding Path): Each decoder block consists of:

\noindent 1. Bilinear Upsampling: Doubles spatial dimensions
\begin{equation}
F_{up}=Upsample_{2x}(F_{enc})
\end{equation}

\noindent 2. Skip Connection: Concatenation with corresponding encoder features
\begin{equation}
F_{concat}=[F_{up};F_{skip}]
\end{equation}

\noindent 3. Double Convolution Block
\begin{equation}
F_{out}=DoubleConv(F_{concat})
\end{equation}
\begin{equation}
DoubleConv(x)=ReLU(BN(Conv_{3\times3}(ReLU(BN(Conv_{3\times3}(x))))))
\end{equation}

\noindent Output Layer: A $1\times1$ convolution maps features to the final segmentation output
\begin{equation}
\tilde{M}=\sigma(Conv_{1\times1}(F_{final}))
\end{equation}
where $\sigma$ is the sigmoid activation function.

\subsubsection{U-Net Loss Function}
Dice-BCE Combined Loss: We employ a combined loss function that addresses both pixel-wise accuracy and region overlap:

\noindent Binary Cross-Entropy Loss:
\begin{equation}
\mathcal{L}_{BCE}=-\frac{1}{N}\sum_{i=1}^{N}[M_{i}log(\tilde{M}_{i})+(1-M_{i})log(1-\tilde{M}_{i})]
\end{equation}

\noindent Dice Loss: The Dice coefficient measures overlap between predicted and ground truth masks:
\begin{equation}
Dice=\frac{2|X\cap Y|}{|X|+|Y|}=\frac{2TP}{2TP+FP+FN}
\end{equation}
\begin{equation}
\mathcal{L}_{Dice}=1-\frac{2\sum_{i=1}^{N}M_{i}\tilde{M}_{i}+\epsilon}{\sum_{i=1}^{N}M_{i}+\sum_{i=1}^{N}\tilde{M}_{i}+\epsilon}
\end{equation}
where $\epsilon$ is a smoothing factor to prevent division by zero.

\noindent Combined Loss:
\begin{equation}
\mathcal{L}_{seg}=\lambda_{BCE}\mathcal{L}_{BCE}+\lambda_{Dice}\mathcal{L}_{Dice}
\end{equation}
Typically, $\lambda_{BCE}=\lambda_{Dice}=0.5$.

\subsubsection{FSM Training Configuration}
\begin{table}[H]
\caption{FSM Training Configuration}
\label{tab:fsm_config}
\centering
\begin{tabular}{ll}
\toprule
\textbf{Hyperparameter} & \textbf{Value} \\
\midrule
Input Size & $256\times256$ \\
Batch Size & 16 \\
Epochs & 50 \\
Optimizer & AdamW \\
Learning Rate & $10^{-4}$ \\
LR Schedule & ReduceLROnPlateau \\
Weight Decay & $10^{-5}$ \\
Encoder Weights & ImageNet pretrained \\
Loss Function & Dice + BCE \\
\bottomrule
\end{tabular}
\end{table}

\subsubsection{Data Augmentation Strategy}
We employed the Albumentations library for on-the-fly augmentation:

\begin{table}[H]
\caption{Augmentation Parameters}
\label{tab:aug_config}
\centering
\begin{tabular}{ll}
\toprule
\textbf{Augmentation} & \textbf{Parameters} \\
\midrule
Horizontal Flip & $p=0.5$ \\
Vertical Flip & $p=0.5$ \\
Random Rotation & limit=$45^{\circ}$, $p=0.5$ \\
Random Brightness/Contrast & brightness=0.2, contrast=0.2, $p=0.3$ \\
Shift Scale Rotate & shift=0.1, scale=0.1, rotate=$15^{\circ}$ \\
Gaussian Noise & var limit=(10,50), $p=0.2$ \\
Grid Distortion & $p=0.2$ \\
\bottomrule
\end{tabular}
\end{table}

\subsubsection{Normalization}
Input images are normalized using ImageNet statistics:
\begin{equation}
x_{norm}=\frac{x-\mu}{\sigma}
\end{equation}
where: $\mu=[0.485,0.456,0.406]$ (RGB channel means) and $\sigma=[0.229,0.224,0.225]$ (RGB channel standard deviations).

\subsection{Complete Inference Pipeline}
The end-to-end inference process operates as follows:

\begin{algorithm}[H]
\caption{DetSEG Inference Pipeline}
\label{alg:inference}
\begin{algorithmic}[1]
\REQUIRE Colonoscopy image I, confidence threshold conf, IoU threshold iou
\ENSURE Detected polyps with bounding boxes and segmentation masks
\STATE I\_resized $\leftarrow$ Resize(I, $640\times640$)
\STATE Detections $\leftarrow$ YOLOv8(I\_resized)
\STATE Detections $\leftarrow$ FilterByConfidence(Detections, conf)
\STATE Detections $\leftarrow$ NMS(Detections, iou)
\STATE M\_final $\leftarrow$ ZeroMatrix(H, W)
\FOR{each detection in Detections}
    \STATE Add contextual padding:
    \STATE $x1 \leftarrow max(0, x-pad), y1 \leftarrow max(0, y-pad)$
    \STATE $x2 \leftarrow min(W, x+pad), y2 \leftarrow min(H, y+pad)$
    \STATE ROI $\leftarrow$ Crop(I, [y1:y2, x1:x2])
    \STATE ROI\_resized $\leftarrow$ Resize(ROI, $256\times256$)
    \STATE ROI\_norm $\leftarrow$ Normalize(ROI\_resized)
    \STATE Logits $\leftarrow$ UNet(ROI\_norm)
    \STATE M\_roi $\leftarrow$ Sigmoid(Logits) $>$ 0.5
    \STATE M\_scaled $\leftarrow$ Resize(M\_roi, original\_dimensions)
    \STATE M\_final $\leftarrow$ BitwiseOR(M\_final, M\_scaled)
\ENDFOR
\RETURN Detections, M\_final
\end{algorithmic}
\end{algorithm}

\subsection{Implementation Details}
The DetSEG framework was implemented using:

\noindent Deep Learning Framework: PyTorch 2.0

\noindent Detection: Ultralytics YOLOv8

\noindent Segmentation: Segmentation Models PyTorch (smp) library

\noindent Data Augmentation: Albumentations

\noindent Image Processing: OpenCV, PIL

\noindent Hardware: NVIDIA GPU with CUDA support

% ============================================================
% 4. EXPERIMENTAL SETUP
% ============================================================
\section{Experimental Setup}

\subsection{Datasets}
\noindent Kvasir-SEG Dataset: 1,000 polyp images with pixel-level annotations - Variable resolution ($332\times487$ to $1920\times1072$) - Diverse polyp sizes, morphologies, and imaging conditions

\noindent CVC-ClinicDB: 612 frames extracted from colonoscopy videos - $384\times288$ resolution - 31 unique polyp sequences

\noindent Combined Dataset Statistics:

\begin{table}[H]
\caption{Combined Dataset Statistics}
\label{tab:dataset_stats}
\centering
\begin{tabular}{lccc}
\toprule
\textbf{Metric} & \textbf{Train} & \textbf{Validation} & \textbf{Test} \\
\midrule
Images & $\sim$1,290 & $\sim$161 & $\sim$161 \\
Split Ratio & 80\% & 10\% & 10\% \\
Polyps/Image & 1-3 & 1-3 & 1-3 \\
\bottomrule
\end{tabular}
\end{table}

\subsection{Evaluation Metrics}

\subsubsection{Detection Metrics}
\noindent Precision (P):
\begin{equation}
P=\frac{TP}{TP+FP}
\end{equation}

\noindent Recall (R):
\begin{equation}
R=\frac{TP}{TP+FN}
\end{equation}

\noindent Average Precision (AP):
\begin{equation}
AP=\int_{0}^{1}P(R)dR
\end{equation}

\noindent Mean Average Precision (mAP@0.5): AP calculated at IoU threshold of 0.5.

\subsubsection{Segmentation Metrics}
\noindent Intersection over Union (IoU / Jaccard Index):
\begin{equation}
IoU=\frac{|M_{pred} \cap M_{gt}|}{|M_{pred} \cup M_{gt}|}=\frac{TP}{TP+FP+FN}
\end{equation}

\noindent Dice Coefficient (F1 Score):
\begin{equation}
Dice=\frac{2|M_{pred} \cap M_{gt}|}{|M_{pred}|+|M_{gt}|}=\frac{2TP}{2TP+FP+FN}
\end{equation}

\noindent Sensitivity (Recall):
\begin{equation}
Sensitivity=\frac{TP}{TP+FN}
\end{equation}

\noindent Specificity:
\begin{equation}
Specificity=\frac{TN}{TN+FP}
\end{equation}

\noindent Precision:
\begin{equation}
Precision=\frac{TP}{TP+FP}
\end{equation}

\subsection{Model Comparison Study}
We conducted a comprehensive comparison study analyzing different object detection architectures for the polyp localization task. Four state-of-the-art detection models were implemented and benchmarked:

\noindent 1. Faster R-CNN: A two-stage detector with region proposal network

\noindent 2. SSD (Single Shot Detector): A single-stage detector with multi-scale feature maps

\noindent 3. RT-DETR: A real-time detection transformer architecture

\noindent 4. YOLOv8: The latest iteration of the YOLO family with anchor-free detection

Selection Criteria: 1. Detection accuracy (mAP@50, mAP@50-95) 2. Precision and Recall 3. Inference speed 4. Robustness to polyp scale variation.

\subsection{Hardware and Training Environment}

\begin{table}[H]
\caption{Hardware and Training Environment}
\label{tab:hardware}
\centering
\begin{tabular}{ll}
\toprule
\textbf{Component} & \textbf{Specification} \\
\midrule
GPU & NVIDIA RTX 3080 / NVIDIA T4 \\
CUDA Version & 11.8 \\
PyTorch Version & 2.0.1 \\
Python Version & 3.10 \\
OS & Windows 11 / Linux \\
Training Time (YOLO) & $\sim$4 hours \\
Training Time (U-Net) & $\sim$2 hours \\
\bottomrule
\end{tabular}
\end{table}

% ============================================================
% 5. RESULTS AND DISCUSSION
% ============================================================
\section{Results and Discussion}

\subsection{Polyp Localization Module Performance}
The YOLOv8s model achieved the following detection performance:

\begin{table}[H]
\caption{PLM Performance}
\label{tab:plm_results}
\centering
\begin{tabular}{ll}
\toprule
\textbf{Metric} & \textbf{Value} \\
\midrule
mAP@0.5 & 0.87 \\
mAP@0.5:0.95 & 0.68 \\
Precision & 0.85 \\
Recall & 0.82 \\
Inference Time & $\sim$15ms \\
\bottomrule
\end{tabular}
\end{table}

Key Observations: High recall ensures minimal polyp miss rate - Anchor-free detection handles irregular polyp shapes effectively - Multi-scale feature fusion captures both small and large polyps.

\subsection{Fine-Grained Segmentation Module Performance}
The U-Net with ResNet-34 encoder achieved:

\begin{table}[H]
\caption{FSM Performance}
\label{tab:fsm_results}
\centering
\begin{tabular}{ll}
\toprule
\textbf{Metric} & \textbf{Value} \\
\midrule
Mean Dice & 0.89 \\
Mean IoU & 0.84 \\
Sensitivity & 0.91 \\
Specificity & 0.98 \\
Precision & 0.87 \\
Inference Time & $\sim$50ms \\
\bottomrule
\end{tabular}
\end{table}

\subsection{Detection Model Comparison}
We compared four state-of-the-art object detection architectures for the Polyp Localization Module.

\begin{table*}[H]
\caption{Detection Model Comparison}
\label{tab:comparison}
\centering
\begin{tabular}{lccccc}
\toprule
\textbf{Model} & \textbf{mAP@50} & \textbf{mAP@50-95} & \textbf{Precision} & \textbf{Recall} & \textbf{F1-Score} \\
\midrule
Faster R-CNN & 92.65\% & 68.51\% & - & 76.65\%* & - \\
SSD & - & - & 83.52\% & 72.73\% & 77.75\% \\
RT-DETR & $\sim$93\% & $\sim$73\% & $\sim$90\% & $\sim$90\% & - \\
YOLOv8 & $\sim$95\% & $\sim$72\% & $\sim$88\% & $\sim$85\% & - \\
\bottomrule
\multicolumn{6}{l}{*Faster R-CNN reports mAR@10 instead of standard recall}
\end{tabular}
\end{table*}

Key Findings:

\noindent 1. YOLOv8 achieved the highest mAP@50 ($\sim$95\%), demonstrating superior detection accuracy for polyp localization

\noindent 2. RT-DETR showed competitive performance with high precision and recall ($\sim$90\% each), but YOLOv8 provided better overall detection rates

\noindent 3. Faster R-CNN achieved 92.65\% mAP@50 with the highest mAP@50-95 (68.51\%), but its two-stage architecture results in slower inference

\noindent 4. SSD showed the lowest overall performance with precision of 83.52\% and recall of 72.73\%

Model Selection Justification: Based on our comparison study, YOLOv8 was selected as the optimal architecture for the Polyp Localization Module due to: Highest mAP@50 detection accuracy ($\sim$95\%) - Single-stage architecture enabling faster inference - Anchor-free design providing better generalization on irregular polyp shapes - Optimal balance between accuracy and computational efficiency for clinical deployment.

% ============================================================
% 6. DISCUSSION
% ============================================================
\section{Discussion}

\subsection{Clinical Applicability}
The DetSEG framework offers several clinical benefits:

\noindent 1. Efficient Inference: Sub-100ms inference enables rapid polyp analysis

\noindent 2. Reduced Variability: Consistent algorithm performance complements endoscopist expertise

\noindent 3. Training Tool: Visual feedback can assist in endoscopy training

\noindent 4. Documentation: Automated detection assists in procedure quality metrics

\subsection{Limitations}

\noindent 1. Dataset Scope: Training limited to publicly available datasets may not represent all clinical scenarios

\noindent 2. Domain Shift: Performance may degrade with equipment from different manufacturers

\noindent 3. Edge Cases: Rare polyp morphologies may be underrepresented

\noindent 4. Integration: Clinical deployment requires regulatory approval and EMR integration

\subsection{Future Directions}

\noindent Video Analysis: Extend the framework to handle colonoscopy video sequences with temporal modeling.

\noindent Uncertainty Quantification: Provide confidence estimates for clinical decision support.

\noindent Polyp Classification: Extend to differentiate adenomatous from hyperplastic polyps.

\noindent Multi-Site Validation: Evaluate generalization across institutions and equipment.

\noindent Lightweight Models: Knowledge distillation for edge deployment.

\noindent Active Learning: Semi-supervised approaches to leverage unlabeled clinical data.

% ============================================================
% 7. CONCLUSION
% ============================================================
\section{Conclusion}

This paper presented DetSEG, a two-stage deep learning framework for automated polyp detection and segmentation in colonoscopy images. By decoupling the localization and segmentation tasks, DetSEG addresses the inherent challenges of polyp scale variation and background complexity. The framework employs YOLOv8 for efficient and accurate polyp detection, followed by a U-Net with ResNet-34 encoder for precise boundary delineation. Experimental results demonstrate that DetSEG achieves competitive performance with state-of-the-art methods while offering unique advantages through its modular design. The system achieves a mean Dice coefficient of 0.891 and IoU of 0.842, with efficient inference times suitable for clinical deployment. The ``detect-then-segment'' approach provides both localization (bounding boxes with confidence scores) and segmentation outputs, offering comprehensive polyp analysis for clinical application. The DetSEG framework represents a significant step toward robust computer-aided detection systems that can assist endoscopists in improving polyp detection rates and ultimately contributing to colorectal cancer prevention.

% ============================================================
% REFERENCES
% ============================================================
\begin{thebibliography}{00}
\bibitem{b1} Siegel, R.L., Miller, K.D., et al. ``Colorectal cancer statistics, 2023'' CA: A Cancer Journal for Clinicians, 2023.
\bibitem{b2} Rex, D.K. et al. ``Colonoscopic miss rates of adenomas determined by back-to-back colonoscopies.'' Gastroenterology, 1997.
\bibitem{b3} Redmon, J., et al. ``You Only Look Once: Unified, Real-Time Object Detection'' CVPR, 2016.
\bibitem{b4} Ronneberger, O., Fischer, P., Brox, T. ``U-Net: Convolutional Networks for Biomedical Image Segmentation'' MICCAI, 2015.
\bibitem{b5} He, K., Zhang, X., et al. ``Deep Residual Learning for Image Recognition'' CVPR, 2016.
\bibitem{b6} Jha, D., et al. ``Kvasir-SEG: A Segmented Polyp Dataset'' MMM, 2020.
\bibitem{b7} Bernal, J., et al. ``WM-DOVA Maps for Accurate Polyp Highlighting in Colonoscopy.'' IEEE Transactions on Medical Imaging, 2015.
\bibitem{b8} Fan, D.P., et al. ``PraNet: Parallel Reverse Attention Network for Polyp Segmentation.'' MICCAI, 2020.
\bibitem{b9} Huang, C.H., et al. ``HarDNet-MSEG: A Simple Encoder-Decoder Polyp Segmentation Neural Network.'' 2021.
\bibitem{b10} Dong, B., et al. ``Polyp-PVT: Polyp Segmentation with Pyramid Vision Transformers.'' CAAI AIR, 2021.
\bibitem{b11} Jocher, G., et al. ``Ultralytics YOLOv8.'' GitHub Repository, 2023.
\bibitem{b12} Yakubovskiy, P. ``Segmentation Models PyTorch'' GitHub Repository, 2019.
\bibitem{b13} Zheng, Z., et al. ``Distance-IoU Loss: Faster and Better Learning for Bounding Box Regression.'' AAAI, 2020.
\bibitem{b14} Lin, T.Y., et al. ``Focal Loss for Dense Object Detection.'' ICCV, 2017.
\bibitem{b15} Zhou, Z., et al. ``UNet++: A Nested U-Net Architecture for Medical Image Segmentation.'' DLMIA, 2018.
\end{thebibliography}

% ============================================================
% APPENDIX
% ============================================================
\appendix

\section{Detailed Network Architectures}
\subsection{YOLOv8s Architecture Summary}
\begin{table}[H]
\caption{YOLOv8s Architecture Summary}
\centering
\begin{tabular}{lcl}
\toprule
\textbf{Layer (type)} & \textbf{Output Shape} & \textbf{Param \#} \\
\midrule
CBS (Conv-BN-SiLU) & [-1, 32, 320, 320] & 464 \\
CBS & [-1, 64, 160, 160] & 18,560 \\
C2f & [-1, 64, 160, 160] & 29,184 \\
CBS & [-1, 128, 80, 80] & 73,984 \\
C2f & [-1, 128, 80, 80] & 197,632 \\
CBS & [-1, 256, 40, 40] & 295,424 \\
C2f & [-1, 256, 40, 40] & 788,992 \\
CBS & [-1, 512, 20, 20] & 1,180,672 \\
C2f & [-1, 512, 20, 20] & 1,838,592 \\
SPPF & [-1, 512, 20, 20] & 656,896 \\
Upsample+Cat+C2f & [-1, 256, 40, 40] & 329,216 \\
Upsample+Cat+C2f & [-1, 128, 80, 80] & 87,040 \\
Concat+C2f & [-1, 256, 40, 40] & 364,544 \\
Concat+C2f & [-1, 512, 20, 20] & 1,313,280 \\
Detect Head & [P3, P4, P5] & $\sim$1.8M \\
\midrule
\textbf{Total params:} & \textbf{$\sim$11.1M} & \\
\bottomrule
\end{tabular}
\end{table}

\subsection{U-Net with ResNet-34 Encoder Summary}
\begin{table}[H]
\caption{U-Net with ResNet-34 Encoder Summary}
\centering
\begin{tabular}{lcl}
\toprule
\textbf{Module} & \textbf{Output Shape} & \textbf{Param \#} \\
\midrule
\textbf{Encoder (ResNet-34):} & & \\
layer0 & [64, 128, 128] & 9,536 \\
layer1 & [64, 64, 64] & 147,968 \\
layer2 & [128, 32, 32] & 525,568 \\
layer3 & [256, 16, 16] & 2,099,712 \\
layer4 & [512, 8, 8] & 8,393,728 \\
\textbf{Decoder:} & & \\
$up1+conv$ & [256, 16, 16] & 3,540,480 \\
$up2+conv$ & [128, 32, 32] & 885,376 \\
$up3+conv$ & [64, 64, 64] & 221,440 \\
$up4+conv$ & [64, 128, 128] & 55,424 \\
\textbf{Output:} & & \\
final conv & [1, 256, 256] & 65 \\
\midrule
\textbf{Total params:} & \textbf{$\sim$24.4M} & \\
\bottomrule
\end{tabular}
\end{table}

\section{Sample Code}
\label{sec:sample_code}

\begin{verbatim}
import torch
import cv2
import numpy as np
from ultralytics import YOLO
import segmentation_models_pytorch as smp
import albumentations as A
from albumentations.pytorch import ToTensorV2

# Load models
yolo_model = YOLO('yolov8s_polyp.pt')
unet_model = smp.Unet(
    encoder_name="resnet34",
    encoder_weights=None,
    in_channels=3,
    classes=1
)
unet_model.load_state_dict(torch.load('unet_best.pth'))
unet_model.eval()
\end{verbatim}

\begin{verbatim}
# Preprocessing
transform = A.Compose([
    A.Resize(256, 256),
    A.Normalize(mean=(0.485, 0.456, 0.406), 
                std=(0.229, 0.224, 0.225)),
    ToTensorV2()
])

def process_image(image, conf_threshold=0.25):
    # Detection
    results = yolo_model(image, conf=conf_threshold)
    h, w = image.shape[:2]
    final_mask = np.zeros((h, w), dtype=np.uint8)
    
    for box in results[0].boxes:
        x1, y1, x2, y2 = box.xyxy[0].cpu().numpy().astype(int)
        
        # Add padding
        pad = 15
        x1, y1 = max(0, x1-pad), max(0, y1-pad)
        x2, y2 = min(w, x2+pad), min(h, y2+pad)
        
        crop_image = image[y1:y2, x1:x2]
        aug = transform(image=crop_image)
        tensor = aug['image'].unsqueeze(0)
        
        with torch.no_grad():
            pred = torch.sigmoid(unet_model(tensor))
            mask = pred[0, 0].cpu().numpy()
            
        mask = (mask > 0.5).astype(np.uint8)
        mask = cv2.resize(mask, (x2-x1, y2-y1))
        
        final_mask[y1:y2, x1:x2] = np.maximum(
            final_mask[y1:y2, x1:x2], mask
        )
        
    return final_mask, results
\end{verbatim}

\section{Hyperparameter Sensitivity Analysis}

\subsection{Learning Rate Impact (U-Net)}
\begin{table}[H]
\caption{Learning Rate Impact}
\centering
\begin{tabular}{lcc}
\toprule
\textbf{Learning Rate} & \textbf{Final Dice} & \textbf{Convergence Epoch} \\
\midrule
$1e-3$ & 0.842 & 35 \\
$5e-4$ & 0.869 & 42 \\
$1e-4$ & 0.891 & 48 \\
$5e-5$ & 0.884 & 55 \\
$1e-5$ & 0.851 & Did not converge \\
\bottomrule
\end{tabular}
\end{table}

\subsection{Batch Size Impact}
\begin{table}[H]
\caption{Batch Size Impact}
\centering
\begin{tabular}{cccc}
\toprule
\textbf{Batch Size} & \textbf{Training Time} & \textbf{Dice} & \textbf{GPU Memory (GB)} \\
\midrule
4 & 4.2h & 0.883 & 3.2 \\
8 & 3.4h & 0.887 & 5.1 \\
16 & 2.8h & 0.891 & 8.7 \\
32 & 2.2h & 0.888 & 15.3 \\
\bottomrule
\end{tabular}
\end{table}

\end{document}